\par\medskip\noindent\textbf{Cómo leer la salida.}
{\small\setlength{\tabcolsep}{4pt}\renewcommand{\arraystretch}{1.15}
\begin{xltabular}{\linewidth}{@{}>{\raggedright\arraybackslash}p{0.24\linewidth}>{\raggedleft\arraybackslash}p{0.15\linewidth}>{\raggedright\arraybackslash}X@{}}
\toprule Indicador & Valor & Interpretación \\ \midrule \endfirsthead
\toprule Indicador & Valor & Interpretación \\ \midrule \endhead
\bottomrule \endlastfoot
\texttt{F(2, 5022)} & 1,870 ($p=0{,}154$) & Prueba $F$ clásica de que los rezagos de $\Delta v$ no entran en la ecuación de $r$; no se rechaza. \\
\texttt{Wald HC0} & 1,947 ($p=0{,}378$) & Misma prueba robusta a heterocedasticidad ($\chi^2$); tampoco rechaza. \\
$r$ no causa a $\Delta v$ ($p=2$) & $F$: $p=0{,}070$; Wald: $p=0{,}142$ & Rechazo marginal con la $F$ clásica que desaparece con la versión robusta. \\
Con $p=7$ & $p$-valores 0,14 a 0,51 & Con más rezagos tampoco hay causalidad de Granger en ninguna dirección. \\
Carteras: orden elegido & AIC 5; BIC 1 & Rezagos del VAR mensual de las carteras. \\
La grande no causa a la pequeña ($p=1$) & $F$: $p=0{,}055$; Wald: $p=0{,}116$ & Evidencia débil de que la grande adelanta a la pequeña (efecto del Capítulo~7); no es significativa con errores robustos. \\
La pequeña no causa a la grande & $p=0{,}56$ y 0,60 & No hay causalidad en esa dirección: la relación, si existe, es asimétrica. \\
Con $p=5$ & $p=0{,}066$ y 0,346 & La conclusión no cambia con más rezagos. \\
\end{xltabular}}
